\phantomsection
\chapter*{Nahkoda di Tengah Badai}
\addcontentsline{toc}{section}{Nahkoda di Tengah Badai --- Umar Prabu Ardiansyah}
\penulis{Umar Prabu Ardiansyah}

\awalcerpen{M}{alam itu}, warna langit berubah menjadi hitam pekat. Suara gemuruh petir bersahut-sahutan di atas Samudra Selatan, membuat dada siapa pun yang mendengarnya merasa ciut. Aku baru dua minggu bergabung di atas kapal layar \(\textit{Bahari Hitam}.\) Bagi seorang anak muda yang baru belajar melaut seperti diriku, malam ini adalah ujian yang sangat menakutkan.

Gelombang air laut naik sangat tinggi, menghantam bagian samping kapal hingga terguncang hebat. Suara kayu lambung kapal yang beradu dengan kerasnya ombak terdengar seperti akan pecah kapan saja. Seluruh \(\textit{kru}\) kapal berlarian ke sana kemari di bawah guyuran hujan deras, berusaha mengikat tali layar agar tidak robek dihantam angin kencang. Aku hanya bisa berdiri terpaku di dekat tiang utama, memegang erat tali tambang dengan tangan yang sudah gemetaran.

``Rian! Jangan cuma berdiri di situ! Bantu tahan tali pelampung!'' teriak Kapten Hendra dari atas \textit{dek} kemudi.

Kapten Hendra adalah ketua kami, seorang pelaut berpengalaman yang sangat disegani oleh seluruh \textit{kru} di kapal. Wajahnya yang penuh guratan pengalaman selalu terlihat tenang, bahkan di tengah cuaca buruk yang membuat seluruh \textit{kru} panik.

Aku langsung berlari mendekati tangga menuju \textit{dek} kemudi, meski beberapa kali hampir tergelincir karena lantai kayu kapal yang sangat licin terkena air hujan dan cipratan air laut.

``Kamu takut, Nak?'' tanya Kapten Hendra saat aku sudah berdiri di sampingnya dengan tubuh yang basah kuyup. Matanya tetap fokus menatap lurus ke arah gelombang laut yang tinggi di depan kami.

``Jujur$\dots$ iya, Kapten. Saya takut kapal kita tidak sanggup bertahan dan tenggelam,'' jawabku jujur dengan suara yang hampir tertelan oleh suara gemuruh angin malam.

Kapten Hendra tidak marah mendengar pengakuanku. Beliau justru memanggilku untuk melangkah lebih dekat ke roda kemudi besar yang terbuat dari kayu jati tua itu. Beliau memegang tanganku, lalu meletakkannya tepat di atas jari-jari roda kemudi yang terasa dingin dan berat.

``Pegang yang kuat,'' kata Kapten Hendra dengan suara berat yang seketika terasa menenangkan. ``Ketakutan di tengah laut itu hal yang sangat wajar. Tapi kalau kamu membiarkan rasa takut itu mengendalikan dirimu, kemudi kapal ini akan terlepas dari genggaman. Saat itulah kita benar-benar berada dalam bahaya.''

Aku merasakan dinginnya kayu kemudi di telapak tanganku, bersanding erat dengan genggaman kuat dari Kapten Hendra. Perlahan, aku menarik napas panjang dan mencoba menenangkan pikiranku. Aku memusatkan seluruh perhatian pada arah haluan kapal, menahan dorongan ombak besar yang terus-menerus mencoba memutar roda kemudi secara mendadak.

Selama hampir dua jam penuh kami bertarung menembus kegelapan malam dan lebatnya hujan. Tenaga kami terkuras habis untuk mempertahankan posisi kapal agar tidak dihantam ombak dari samping. Sedikit demi sedikit, ketakutan yang sejak awal melumpuhkan diriku perlahan menghilang, berganti menjadi rasa fokus dan keberanian yang muncul dari dalam diri.

Saat fajar perlahan datang dari ufuk timur, awan hitam pekat mulai terangkat. Sinar matahari pagi menerobos dari balik awan, menyinari permukaan laut yang mendadak berubah menjadi tenang dan membiru. Kapal kami berhasil melewati area badai tanpa ada kerusakan parah pada bagian layar maupun lambung.

Aku terduduk lemas di atas geladak yang basah, menatap telapak tanganku yang memerah dan lecet akibat menahan tali dan kemudi semalaman. Kapten Hendra berjalan mendekat ke arahku, lalu menepuk bahuku pelan sambil tersenyum tipis.

``Kamu sudah melakukannya dengan sangat baik hari ini,'' ujar Kapten Hendra singkat sebelum beliau berjalan kembali menuju kabinnya.

Pagi itu aku belajar satu hal yang sangat berharga di atas kapal. Keberanian itu bukan berarti kita tidak punya rasa takut sama sekali, melainkan bagaimana kita mampu mengendalikan rasa takut tersebut agar tetap bisa melangkah dan menjaga kemudi kehidupan kita sendiri.

\jedahias
